% SolarMax — Wiring Reference (typeset companion to wiring.md)
% Compile:  pdflatex wiring.tex   (run from docs/)
% This is the human-readable, print-friendly version of wiring.md. The tables
% here are the same authoritative per-pin detail, formatted for easy reading.

\documentclass[11pt]{article}
\usepackage[T1]{fontenc}
\usepackage[margin=1.8cm]{geometry}
\usepackage{lmodern}
\usepackage[table,svgnames]{xcolor}
\usepackage{booktabs}
\usepackage{array}
\usepackage{tabularx}
\usepackage{xltabular}
\usepackage{ragged2e}
\usepackage{amssymb}
\usepackage{textcomp}
\usepackage{enumitem}
\usepackage[most]{tcolorbox}
\usepackage{hyperref}
\usepackage{underscore}   % lets us write _ literally in pin/file names

\hypersetup{colorlinks=true, linkcolor=NavyBlue, urlcolor=NavyBlue}

% ---- palette ----------------------------------------------------------------
\definecolor{headblue}{HTML}{1F4E79}
\definecolor{rowgray}{HTML}{EEF3F8}
\definecolor{okgreen}{HTML}{2E7D32}
\definecolor{warnorange}{HTML}{B34700}
\definecolor{infoblue}{HTML}{1565C0}

% ---- status marks -----------------------------------------------------------
\newcommand{\yes}{\textcolor{okgreen}{\textbf{\checkmark\ confirmed}}}
\newcommand{\caution}{\textcolor{warnorange}{\textbf{!\ }}}
\newcommand{\info}{\textcolor{infoblue}{\textbf{i\ }}}

% ---- table column types -----------------------------------------------------
\newcolumntype{P}[1]{>{\RaggedRight\arraybackslash}p{#1}}
\newcolumntype{C}[1]{>{\Centering\arraybackslash}p{#1}}
\newcolumntype{T}[1]{>{\RaggedRight\arraybackslash\ttfamily\footnotesize}p{#1}}
\newcolumntype{L}{>{\RaggedRight\arraybackslash}X}
\newcolumntype{N}{>{\Centering\arraybackslash}p{0.5cm}}
\newcolumntype{Q}{>{\Centering\arraybackslash}p{0.85cm}}

\renewcommand{\arraystretch}{1.25}
\newcommand{\tc}[1]{\textcolor{white}{\textbf{#1}}}   % header cell text
\newcommand{\nd}{\textrm{\textemdash}}                % "not applicable" cell (roman em-dash, ligature-safe in tt cols)

% ---- callout boxes ----------------------------------------------------------
\newtcolorbox{warnbox}{breakable, enhanced, colback=warnorange!6,
  colframe=warnorange!75, boxrule=0.6pt, arc=2pt,
  left=7pt, right=7pt, top=5pt, bottom=5pt}
\newtcolorbox{notebox}{breakable, enhanced, colback=infoblue!5,
  colframe=infoblue!55, boxrule=0.5pt, arc=2pt,
  left=7pt, right=7pt, top=5pt, bottom=5pt}

% tighter lists inside boxes
\setlist[enumerate]{leftmargin=1.6em, topsep=2pt, itemsep=2pt}
\setlist[itemize]{leftmargin=1.4em, topsep=2pt, itemsep=2pt}

\newcommand{\sectionrule}{\vspace{2pt}\hrule height 0.4pt\vspace{8pt}}

\begin{document}

% =============================================================================
\begin{center}
  {\LARGE\bfseries SolarMax \nd\ Wiring Reference}\\[4pt]
  {\small All pin assignments are defined in \texttt{include/config.h}.
   The tables below are the authoritative per-pin detail; use them with the
   colored diagrams in \texttt{wiring\_diagram.pdf} and
   \texttt{wiring\_breadboard.pdf}.}
\end{center}
\vspace{4pt}

\begin{notebox}
\textbf{Related files}
\begin{itemize}
  \item \texttt{breadboard\_placement.md} \nd\ column-by-column placement map
        (single source of truth for the breadboard).
  \item \texttt{wiring\_breadboard.pdf} \nd\ \textbf{build from this.} Full
        breadboard layout, every wire drawn.
  \item \texttt{wiring\_diagram.pdf} \nd\ schematic overview (power distribution,
        signal map, voltage-divider details). Good for understanding
        \emph{what connects to what}.
\end{itemize}
Rebuild any diagram with \texttt{pdflatex <file>.tex}.
\end{notebox}

% =============================================================================
\section*{Datasheet verification status}
Checked against the PDFs in \texttt{../datasheets/} on 2026-06-19.

\smallskip
\rowcolors{2}{rowgray}{white}
\begin{xltabular}{\textwidth}{@{}L T{3.2cm} L@{}}
\toprule
\rowcolor{headblue}\tc{Part} & \tc{Datasheet} & \tc{Result}\\
\midrule\endhead
ESP32 pin map (input-only 34--39, ADC1, no strapping pins used)
  & esp32_datasheet_en.pdf & \yes\\
BTS7960 pinout + 3.3\,V logic + 6--27\,V motor supply
  & BTS7960 Motor Driver.pdf & \yes\ (resolved the level-shifter question)\\
DS3231 I2C address 0x68 & DS3231.pdf & \yes\\
Potentiometer (P160 rotary, 3-terminal) & p160.pdf & \yes\\
\textbf{Anemometer = QS-FS powered sensor, not a reed switch}
  & C2192-002_datasheet.pdf
  & \caution\textbf{wiring corrected} \nd\ see anemometer section\\
LDR / photocell & SEN-09088.pdf
  & \caution image-only PDF, not machine-readable; wiring is standard
    2-terminal photocell\\
Solar panel (Renogy 50\,W) & 50D-SS-Datasheet.pdf
  & \info this is the tracked panel, not control wiring\\
\bottomrule
\end{xltabular}

\smallskip
\noindent\textbf{Motor:} 5882-50ZY identified from the model number
(stall $\sim$15--19\,A, max 200\,kg$\cdot$cm). \textbf{Confirm the gear-ratio/RPM
variant ordered} \nd\ only the high-ratio (low-RPM) versions meet the 20\,N$\cdot$m
wind-holding requirement; the fast 440\,RPM version is $\sim$10$\times$ too weak.

\smallskip
\noindent\textbf{Limit switches:} REIFENG mechanical endstop selected (SPDT
C/NO/NC, 1\,A 125\,VAC). Generic 3D-printer part with no formal datasheet ---
specs read directly from the switch body.

% =============================================================================
\section*{Components List}
\rowcolors{2}{rowgray}{white}
\begin{xltabular}{\textwidth}{@{}N P{3.6cm} Q >{\RaggedRight\arraybackslash\ttfamily\scriptsize}p{2.8cm} L@{}}
\toprule
\rowcolor{headblue}\tc{\#} & \tc{Component} & \tc{Qty} & \tc{Datasheet} & \tc{Notes}\\
\midrule\endhead
1 & ESP32 DevKit V1 (WROOM-32) & 1
  & esp32_datasheet_en.pdf & 3.3\,V logic (also esp-dev-kits-en-master-esp32.pdf)\\
2 & BTS7960 43A H-Bridge module & 1 & BTS7960 Motor Driver.pdf
  & Motor driver, 3.3--5\,V logic, 6--27\,V motor supply\\
3 & DC worm gear motor, 12\,V (5882-50ZY) & 1 & vendor spec (model identified)
  & Brushed, self-locking; stall $\sim$15--19\,A; \textbf{confirm gear-ratio
    variant} \nd\ needs $\sim$200\,kg$\cdot$cm / 20\,N$\cdot$m for wind\\
4 & P160 16\,mm rotary potentiometer & 1 & p160.pdf
  & Use \textbf{linear (B) taper}; 10\,k$\Omega$; 260\textdegree\ travel\\
5 & DS3231 RTC module & 1 & DS3231.pdf & I2C 0x68, CR2032 backup\\
6 & REIFENG mechanical endstop (SPDT microswitch on PCB, C/NO/NC,
    1\,A 125\,VAC, 1\,M cable) & 2 & generic 3D-printer part (no datasheet)
  & East \& west end-stops; use C + NO contacts\\
7 & QS-FS wind speed sensor (C2192-002) & 1 & C2192-002_datasheet.pdf
  & \textbf{Powered 5--24\,V sensor, NOT a reed switch} \nd\ see anemometer
    section\\
8 & LDR / photocell (SparkFun SEN-09088) & 2 & SEN-09088.pdf
  & Diagnostic only; 2-terminal, no polarity\\
9 & 10\,k$\Omega$ resistor & 2 & \nd\ & LDR voltage divider pull-downs\\
10 & 10\,k$\Omega$ + 20\,k$\Omega$ resistor & 1 ea & \nd
  & Anemometer 5\,V$\rightarrow$3.3\,V signal divider (see note)\\
11 & 12\,V 20\,A power supply (e.g.\ ALITOVE 240\,W) & 1 & \nd
  & Motor rail; 20\,A covers the $\sim$19\,A stall surge\\
12 & LM2596 buck converter (12\,V$\rightarrow$5\,V, adj., 3\,A) & 1 & \nd
  & Makes the 5\,V rail; \textbf{set to 5.0\,V before connecting ESP32}\\
13 & 12\,AWG wire & \nd\ & \nd
  & Motor power path only (supply$\leftrightarrow$driver$\leftrightarrow$motor)\\
14 & 18--22\,AWG wire & \nd\ & \nd\ & All logic/signal/ESP32 wiring\\
15 & Schottky diode (1N5822 3\,A, or 1N5817) & 1 & \nd
  & In buck 5\,V output $\rightarrow$ ESP 5V pin (stripe toward ESP); lets USB
    stay safely pluggable for programming\\
--- & Renogy 50\,W solar panel (RNG-50D-SS) & 1 & 50D-SS-Datasheet.pdf
  & The tracked panel (Voc 21.8\,V, Isc 3.1\,A); not wired to the controller\\
\bottomrule
\end{xltabular}

\begin{notebox}
\textbf{Common ground:} ESP32 GND, BTS7960 logic GND, anemometer ground, and the
12\,V supply GND must all be tied together.
\end{notebox}

% =============================================================================
\section*{Power \nd\ two rails}
A single 12\,V 20\,A supply feeds the motor directly and an LM2596 buck converter
that produces the 5\,V logic rail.

\begin{tcolorbox}[colback=gray!6, colframe=gray!40, boxrule=0.4pt, arc=2pt,
  fontupper=\ttfamily\scriptsize]
\begin{verbatim}
12V 20A supply
  |-- B+/B- ------------> BTS7960 --> motor (M+/M-)     [12 AWG, ~19 A]
  |-- buck --[Schottky]-> 5V rail --+-- ESP32 5V pin
                                    |-- BTS7960 VCC
                                    +-- anemometer power
USB port = programming only (the diode lets you plug it in anytime)
ESP32 3.3V --> RTC, pot, LDRs, limit switches          [18-22 AWG]
ALL grounds tied together (12V supply, buck, ESP32, driver, sensors)
\end{verbatim}
\end{tcolorbox}

\smallskip
\rowcolors{2}{rowgray}{white}
\begin{xltabular}{\textwidth}{@{}P{3.2cm} P{4.4cm} P{2.1cm} L@{}}
\toprule
\rowcolor{headblue}\tc{From} & \tc{To} & \tc{Wire} & \tc{Notes}\\
\midrule\endhead
12\,V supply (+) & BTS7960 B+ & 12\,AWG & Motor power rail\\
12\,V supply ($-$) & BTS7960 B$-$ & 12\,AWG & Common GND\\
12\,V supply (+/$-$) & LM2596 IN+ / IN$-$ & 18--22\,AWG & Buck converter input\\
Buck OUT+ & \textbf{Schottky diode} $\rightarrow$ ESP32 5V pin, BTS7960 VCC,
  anemometer power & 18--22\,AWG & 5\,V logic rail; diode lets USB coexist\\
Buck OUT$-$ & ESP32 GND (common) & 18--22\,AWG & \\
ESP32 3.3\,V & Potentiometer pin 1 & 22\,AWG & Pot reference voltage\\
ESP32 3.3\,V & DS3231 VCC & 22\,AWG & RTC power (3.3\,V \nd\ see RTC safety note)\\
ESP32 GND & All peripheral GNDs & 22\,AWG & Common GND \nd\ must be shared\\
\bottomrule
\end{xltabular}

\begin{warnbox}
\begin{itemize}
  \item If using an adjustable buck, \textbf{set it to 5.0\,V before wiring it to
        the ESP32}.
  \item Do \textbf{not} feed 12\,V into the ESP32 VIN; power it from the regulated
        5\,V rail.
\end{itemize}
\end{warnbox}

\begin{notebox}
\textbf{Programming:} the ESP32 is powered through its \textbf{5V pin} from the
buck \nd\ the USB port is for flashing code only. A \textbf{Schottky diode} in the
buck's 5\,V output means you can leave the whole system powered and plug a USB
cable into the computer anytime, with any cable, with no conflict and no risk to
the USB port (the diode blocks any backfeed toward the computer).
\end{notebox}

% =============================================================================
\section*{BTS7960 Motor Driver $\rightarrow$ ESP32}
\rowcolors{2}{rowgray}{white}
\begin{xltabular}{\textwidth}{@{}P{2.2cm} P{2.4cm} T{3cm} L@{}}
\toprule
\rowcolor{headblue}\tc{BTS7960 Pin} & \tc{ESP32 GPIO} & \tc{Config symbol} & \tc{Description}\\
\midrule\endhead
RPWM & GPIO 25 & PIN_MOTOR_RPWM & PWM \nd\ drives motor CW (panel tilts west)\\
LPWM & GPIO 26 & PIN_MOTOR_LPWM & PWM \nd\ drives motor CCW (panel tilts east)\\
R_EN & GPIO 27 & PIN_MOTOR_REN & Enable right half-bridge (driven HIGH)\\
L_EN & GPIO 14 & PIN_MOTOR_LEN & Enable left half-bridge (driven HIGH)\\
VCC & ESP32 5\,V & \nd\ & Logic supply for the module\\
GND & ESP32 GND & \nd\ & Logic ground\\
R_IS & \nd\ & \nd\ & Current-sense output, not connected\\
L_IS & \nd\ & \nd\ & Current-sense output, not connected\\
\bottomrule
\end{xltabular}

\begin{notebox}
\textbf{3.3\,V logic confirmed by datasheet.} The BTS7960 module datasheet states
\textbf{``Control Input Level: 3.3\textasciitilde5V''} \nd\ the ESP32's 3.3\,V GPIO
drives RPWM/LPWM/R\_EN/L\_EN directly, no level shifter needed. Motor supply
(B+/B$-$) range is \textbf{6--27\,VDC}, VCC is +5\,V.
\end{notebox}

% =============================================================================
\section*{BTS7960 $\rightarrow$ Motor \& Power Terminals}
Screw-terminal block labels on the IBT-2 module (verified against
datasheet/vendor docs).

\smallskip
\rowcolors{2}{rowgray}{white}
\begin{xltabular}{\textwidth}{@{}P{2.2cm} P{4.4cm} L@{}}
\toprule
\rowcolor{headblue}\tc{Terminal} & \tc{Connection} & \tc{Notes}\\
\midrule\endhead
M+ & Motor wire A & Swap A/B if the panel drives the wrong direction\\
M$-$ & Motor wire B & \\
B+ & 12\,V supply (+) & Motor power rail, 6--27\,V\\
B$-$ & 12\,V supply ($-$) & Common ground\\
\bottomrule
\end{xltabular}

\begin{notebox}
\textbf{Power sizing:} the 5882-50ZY draws \textbf{$\sim$15--19\,A at stall}
(startup inrush or a jam), though normal slow tracking moves draw far less. The
BTS7960 (43\,A) handles this easily, but the \textbf{12\,V supply/battery needs
$\ge$15--20\,A surge capability} \nd\ a 10\,A supply may brown out on stall.
Fuse $\sim$10--15\,A on B+; 18\,AWG wire is near its limit at stall current.
\end{notebox}

% =============================================================================
\section*{Potentiometer (Shaft Angle Feedback)}
\rowcolors{2}{rowgray}{white}
\begin{xltabular}{\textwidth}{@{}P{3cm} P{3.6cm} L@{}}
\toprule
\rowcolor{headblue}\tc{Pot pin} & \tc{Connection} & \tc{Notes}\\
\midrule\endhead
Pin 1 (end) & ESP32 3.3\,V & Reference high\\
Pin 2 (wiper) & ESP32 GPIO 34 & \texttt{PIN_POT} \nd\ ADC1\_CH6\\
Pin 3 (end) & ESP32 GND & Reference low\\
\bottomrule
\end{xltabular}

\begin{notebox}
Use ADC1 pins only (GPIO 32--39). ADC2 is disabled while WiFi is active.
Run \texttt{runCalibration()} on first assembly to find \texttt{POT_ADC_MIN} /
\texttt{POT_ADC_MAX}.
\end{notebox}

% =============================================================================
\section*{Limit Switches \nd\ REIFENG mechanical endstop (SPDT, C/NO/NC)}
One switch at each end of travel. The board has an SPDT microswitch (terminals
\textbf{C / NO / NC}, rated 1\,A 125\,VAC) on a small PCB with mounting holes and
a 3-wire cable (red/black/white). Firmware uses \texttt{INPUT_PULLDOWN} +
\texttt{LIMIT_ACTIVE HIGH}, so the switch must connect the GPIO to \textbf{3.3\,V}
when pressed \nd\ use the \textbf{C and NO} contacts.

\smallskip
\rowcolors{2}{rowgray}{white}
\begin{xltabular}{\textwidth}{@{}P{3cm} P{2.4cm} T{2.8cm} L@{}}
\toprule
\rowcolor{headblue}\tc{Switch} & \tc{ESP32 GPIO} & \tc{Config symbol} & \tc{Triggers when\dots}\\
\midrule\endhead
CW (west) limit & GPIO 32 & PIN_LIMIT_CW & Panel reaches full west travel\\
CCW (east) limit & GPIO 33 & PIN_LIMIT_CCW & Panel reaches full east travel\\
\bottomrule
\end{xltabular}

\smallskip
\rowcolors{2}{rowgray}{white}
\begin{xltabular}{\textwidth}{@{}P{4cm} L@{}}
\toprule
\rowcolor{headblue}\tc{Switch contact} & \tc{Connect to}\\
\midrule\endhead
C (common) & ESP32 3.3\,V\\
NO (normally open) & ESP32 GPIO 32 (CW) or 33 (CCW)\\
NC (normally closed) & not connected\\
\bottomrule
\end{xltabular}

\begin{notebox}
\textbf{Identify the wires} (red/black/white order isn't guaranteed) with a
multimeter in continuity mode, or with the bring-up tool:
\begin{enumerate}
  \item Lever \textbf{released} $\rightarrow$ the two wires that beep are
        \textbf{C and NC}.
  \item Lever \textbf{pressed} $\rightarrow$ the two wires that beep are
        \textbf{C and NO}.
  \item The wire common to both is \textbf{C} ($\rightarrow$ 3.3\,V); the other
        from step 2 is \textbf{NO} ($\rightarrow$ GPIO).
\end{enumerate}
Not weather-sealed \nd\ mount inside the electronics enclosure for a permanent
outdoor install. (To wire C+NC instead, change \texttt{LIMIT_ACTIVE} to
\texttt{LOW} and the pinMode to \texttt{INPUT_PULLUP}.)
\end{notebox}

% =============================================================================
\section*{Anemometer (Wind Speed) \nd\ QS-FS, part C2192-002}
\begin{warnbox}
\caution\textbf{This is NOT a passive reed switch.} Per
\texttt{datasheets/C2192-002\_datasheet.pdf}, the QS-FS is a \textbf{powered
electromagnetic-induction sensor} that needs its own supply and outputs an
actively-driven signal. This changes the wiring substantially from the earlier
reed-switch assumption.
\end{warnbox}

\noindent\textbf{Specs:} Supply 5--24\,VDC $\cdot$ output options
(variant-dependent): \textbf{Pulse (5\,V)}, 4--20\,mA, 0--5\,V/1--5\,V, or RS485
$\cdot$ range 0--60\,m/s $\cdot$ starting threshold $\le$ 0.8\,m/s $\cdot$ IP55.

\smallskip
\noindent\textbf{3-wire connection (Pulse / current / voltage variants):}

\smallskip
\rowcolors{2}{rowgray}{white}
\begin{xltabular}{\textwidth}{@{}P{2.2cm} P{4cm} L@{}}
\toprule
\rowcolor{headblue}\tc{Function} & \tc{Wire color (2 schemes)} & \tc{Connect to}\\
\midrule\endhead
Power & red \textbf{or} brown & ESP32 5\,V (VIN) \nd\ sensor needs 5--24\,V, use 5\,V\\
Ground & blue \textbf{or} black & ESP32 GND (common)\\
Signal & yellow \textbf{or} blue & GPIO 35 (\texttt{PIN_ANEMOMETER})
  \textbf{via level shift \nd\ see below}\\
\bottomrule
\end{xltabular}

\begin{warnbox}
\textbf{Two things to resolve before wiring this:}
\begin{enumerate}
  \item \textbf{Confirm which output variant you bought.} The current firmware
        counts digital pulses (\texttt{RISING}-edge interrupt), so it only works
        with the \textbf{Pulse (5\,V)} variant. If you have the 4--20\,mA, 0--5\,V,
        or RS485 version, the firmware and wiring must change (ADC + sense
        resistor, or a MAX485/UART transceiver). The ``-002'' suffix likely
        encodes the variant \nd\ verify with the supplier.
  \item \textbf{5\,V signal must be dropped to 3.3\,V.} The pulse output is
        \textbf{5\,V}, but GPIO 35's absolute max is $\sim$3.6\,V \nd\ driving it
        at 5\,V can damage the pin. Use a \textbf{voltage divider} on the signal
        line: signal $\rightarrow$ \textbf{10\,k$\Omega$} $\rightarrow$ GPIO 35,
        and GPIO 35 $\rightarrow$ \textbf{20\,k$\Omega$} $\rightarrow$ GND (gives
        5\,V $\times$ 20/30 $\approx$ 3.3\,V). A logic-level shifter also works.
  \item \textbf{Recalibrate} \texttt{ANEM_MPH_PER_HZ}. The default \texttt{1.49}
        was for a SparkFun reed-switch anemometer and is wrong for the QS-FS. Get
        the pulse-frequency-to-wind-speed relation from the QS-FS manual/supplier
        and update \texttt{config.h}.
\end{enumerate}
\end{warnbox}

% =============================================================================
\section*{LDR Sensors (Diagnostic)}
Voltage divider: 3.3\,V $\rightarrow$ LDR $\rightarrow$ ADC pin $\rightarrow$
10\,k$\Omega$ to GND. Used only to cross-check the astronomical algorithm \nd\ not
required for tracking.

\smallskip
\rowcolors{2}{rowgray}{white}
\begin{xltabular}{\textwidth}{@{}P{2.4cm} P{2.4cm} T{2.8cm} L@{}}
\toprule
\rowcolor{headblue}\tc{LDR} & \tc{ESP32 GPIO} & \tc{Config symbol} & \tc{Position}\\
\midrule\endhead
East LDR & GPIO 36 & PIN_LDR_EAST & Faces east (morning sun)\\
West LDR & GPIO 39 & PIN_LDR_WEST & Faces west (afternoon sun)\\
\bottomrule
\end{xltabular}

\smallskip
\rowcolors{2}{rowgray}{white}
\begin{xltabular}{\textwidth}{@{}P{6cm} L@{}}
\toprule
\rowcolor{headblue}\tc{Node} & \tc{Connection}\\
\midrule\endhead
LDR top & ESP32 3.3\,V\\
LDR bottom / ADC junction & ESP32 GPIO 36 or 39\\
10\,k$\Omega$ resistor (top) & ADC junction\\
10\,k$\Omega$ resistor (bottom) & ESP32 GND\\
\bottomrule
\end{xltabular}

% =============================================================================
\section*{DS3231 RTC Module (I2C)}
\rowcolors{2}{rowgray}{white}
\begin{xltabular}{\textwidth}{@{}P{2.6cm} P{2.4cm} T{2.6cm} L@{}}
\toprule
\rowcolor{headblue}\tc{DS3231 pin} & \tc{ESP32 GPIO} & \tc{Config symbol} & \tc{Notes}\\
\midrule\endhead
VCC & \textbf{3.3\,V} & \nd\ & Power from 3.3\,V \nd\ see safety note below\\
GND & GND & \nd\ & \\
SDA & GPIO 21 & RTC_SDA & I2C data, address 0x68\\
SCL & GPIO 22 & RTC_SCL & I2C clock\\
SQW / INT & \nd\ & \nd\ & Not connected\\
32K & \nd\ & \nd\ & Not connected\\
\bottomrule
\end{xltabular}

\begin{warnbox}
\textbf{Power from 3.3\,V, not 5\,V (safety).} The common ZS-042 DS3231 board has
an onboard battery-charging circuit (diode + 200\,$\Omega$) intended for a
rechargeable LIR2032. If you power VCC at 5\,V while using a \textbf{non-rechargeable
CR2032}, it pushes $\sim$5\,mA into a cell that cannot take a charge \nd\ risk of
overheating/leakage. Powering from 3.3\,V keeps the diode reverse-biased and
prevents charging. (If you must run 5\,V, fit a LIR2032 or remove the charging
resistor.)
\end{warnbox}

\begin{notebox}
\begin{itemize}
  \item The DS3231 chip's I2C address is \textbf{0x68}; the onboard EEPROM is 0x57
        (or 0x50--0x57). The bring-up I2C scan should report 0x68.
  \item Install a CR2032 coin cell so the RTC keeps time through power loss.
  \item The module includes onboard I2C pull-up resistors (typically
        4.7\,k$\Omega$) \nd\ do not add external pull-ups.
\end{itemize}
\end{notebox}

% =============================================================================
\section*{ESP32 GPIO Summary}
\rowcolors{2}{rowgray}{white}
\begin{xltabular}{\textwidth}{@{}C{1.4cm} P{3.4cm} P{3.2cm} L@{}}
\toprule
\rowcolor{headblue}\tc{GPIO} & \tc{Function} & \tc{Direction} & \tc{Notes}\\
\midrule\endhead
14 & Motor L_EN & Output & BTS7960 left enable\\
21 & I2C SDA & Bidirectional & DS3231 RTC\\
22 & I2C SCL & Output & DS3231 RTC\\
25 & Motor RPWM & Output & PWM CW\\
26 & Motor LPWM & Output & PWM CCW\\
27 & Motor R_EN & Output & BTS7960 right enable\\
32 & Limit CW & Input (pull-down) & West end-stop\\
33 & Limit CCW & Input (pull-down) & East end-stop\\
34 & Potentiometer & ADC1 input & Shaft angle feedback\\
35 & Anemometer & Input-only & QS-FS pulse \nd\ \textbf{5\,V signal, must divide
  to 3.3\,V} (10\,k/20\,k)\\
36 & LDR East & ADC1 input & Diagnostic\\
39 & LDR West & ADC1 input & Diagnostic\\
\bottomrule
\end{xltabular}

\end{document}
